%% 
%% Copyright 2007-2024 Elsevier Ltd
%% 
%% This file is part of the 'Elsarticle Bundle'.
%% ---------------------------------------------
%% 
%% It may be distributed under the conditions of the LaTeX Project Public
%% License, either version 1.3 of this license or (at your option) any
%% later version.  The latest version of this license is in
%%    http://www.latex-project.org/lppl.txt
%% and version 1.3 or later is part of all distributions of LaTeX
%% version 1999/12/01 or later.
%% 
%% The list of all files belonging to the 'Elsarticle Bundle' is
%% given in the file `manifest.txt'.
%% 
%% Template article for Elsevier's document class `elsarticle'
%% with harvard style bibliographic references

\documentclass[preprint,12pt,authoryear]{elsarticle}

%% Use the option review to obtain double line spacing
%% \documentclass[authoryear,preprint,review,12pt]{elsarticle}

%% Use the options 1p,twocolumn; 3p; 3p,twocolumn; 5p; or 5p,twocolumn
%% for a journal layout:
%% \documentclass[final,1p,times,authoryear]{elsarticle}
%% \documentclass[final,1p,times,twocolumn,authoryear]{elsarticle}
%% \documentclass[final,3p,times,authoryear]{elsarticle}
%% \documentclass[final,3p,times,twocolumn,authoryear]{elsarticle}
%% \documentclass[final,5p,times,authoryear]{elsarticle}
%% \documentclass[final,5p,times,twocolumn,authoryear]{elsarticle}

%% For including figures, graphicx.sty has been loaded in
%% elsarticle.cls. If you prefer to use the old commands
%% please give \usepackage{epsfig}

%% The amssymb package provides various useful mathematical symbols
\usepackage{amssymb}
%% The amsmath package provides various useful equation environments.
\usepackage{amsmath}
%% The amsthm package provides extended theorem environments
%% \usepackage{amsthm}

%% The lineno packages adds line numbers. Start line numbering with
%% \begin{linenumbers}, end it with \end{linenumbers}. Or switch it on
%% for the whole article with \linenumbers.
%% \usepackage{lineno}

\usepackage{amssymb}
\usepackage{xspace}

\newcommand{\even}          {\raisebox{0.3ex}{\scalebox{0.7}{$\bigcirc$}}\xspace}
\newcommand{\odd}           {\scalebox{0.9}{$\square$}\xspace}
\newcommand{\subeven}       {\raisebox{0.3ex}{\scalebox{0.5}{$\bigcirc$}}\xspace}
\newcommand{\subodd}        {\scalebox{0.6}{$\square$}\xspace}
\newcommand{\player}        {\mathsf{p}\xspace}
\newcommand{\opponent}      {\overline{\mathsf{p}}\xspace}

\newcommand{\wregularform}  {$\omega$-regular form\xspace}
\newcommand{\wregulargames} {$\omega$-regular games\xspace}
\newcommand{\noc}           {{\small\textsf{No\textit{-}Opponent\textit{-}Cycle}}\xspace}
\newcommand{\nocchecker}    {{\small\textsf{NOC\textit{-}Checker}}\xspace}
\newcommand{\noceager}      {{\small\textsf{NOC\textit{-}Eager}}\xspace}
\newcommand{\nocfilter}     {{\small\textsf{NOC\textit{-}Filter}}\xspace}
\newcommand{\nocmemo}       {{\small\textsf{NOC\textit{-}Memo}}\xspace}
\newcommand{\nocq}          {{\small\textsf{NOCQ}}\xspace}
\newcommand{\zra}           {{\small\textsf{ZRA}}\xspace}

% \newtheorem{definition}{Definition}
% \newtheorem{example}{Example}

% % Define MiniZinc syntax highlighting
% \lstdefinelanguage{MiniZinc}{
%   keywords={include, constraint, output, array, of, var, bool},
%   morekeywords={[2] set, int, float, string},
%   keywordstyle=\color{blue}\bfseries,
%   keywordstyle={[2]\color{teal}\bfseries},
%   sensitive=true,
%   comment=[l]{\%}, 
%   morecomment=[s]{/*}{*/},
%   commentstyle=\color{gray}\ttfamily,
%   stringstyle=\color{red}\ttfamily,
%   morestring=[b]",    % Defines double-quoted strings
%   morestring=[b]',    % Defines single-quoted strings  basicstyle=\ttfamily\small,
%   numberstyle=\tiny\color{gray},
%   numbers=left,
%   stepnumber=1,
%   breaklines=true,
%   showstringspaces=false,
%   frame=single,
%   tabsize=2,
%   basicstyle=\ttfamily\footnotesize
% }

% \lstset{
%     basicstyle=\small\ttfamily,
%     backgroundcolor=\color{gray!10},
%     frame=single,
%     framerule=0pt,
%     columns=fullflexible,
%     breaklines=true,
%     xleftmargin=0.5em,
%     xrightmargin=0.5em
% }

\journal{Artificial Intelligence}

\begin{document}

\begin{frontmatter}

%% Title, authors and addresses

%% use the tnoteref command within \title for footnotes;
%% use the tnotetext command for theassociated footnote;
%% use the fnref command within \author or \affiliation for footnotes;
%% use the fntext command for theassociated footnote;
%% use the corref command within \author for corresponding author footnotes;
%% use the cortext command for theassociated footnote;
%% use the ead command for the email address,
%% and the form \ead[url] for the home page:
%% \title{Title\tnoteref{label1}}
%% \tnotetext[label1]{}
%% \author{Name\corref{cor1}\fnref{label2}}
%% \ead{email address}
%% \ead[url]{home page}
%% \fntext[label2]{}
%% \cortext[cor1]{}
%% \affiliation{organization={},
%%            addressline={}, 
%%            city={},
%%            postcode={}, 
%%            state={},
%%            country={}}
%% \fntext[label3]{}

\title{Constraint base approach for solving deterministic \wregulargames} %% Article title

%% use optional labels to link authors explicitly to addresses:
\author[label1,label3]{Gonzalo Hernandez}
\author[label1]{Julian Garcia}
\author[label2]{Julian Gutierrez}
\author[label3]{Guido Tack}

\affiliation[label1]{organization={Monash University},
  addressline={},
  city={},
  postcode={},
  state={},
  country={Australia}}

\affiliation[label2]{organization={University of Sussex},
  addressline={},
  city={},
  postcode={},
  state={},
  country={United Kingdom}}

\affiliation[label2]{organization={Universidad de Nariño},
  addressline={},
  city={},
  postcode={},
  state={},
  country={Colombia}}

% \author{} %% Author name

% %% Author affiliation
% \affiliation{organization={},%Department and Organization
%             addressline={}, 
%             city={},
%             postcode={}, 
%             state={},
%             country={}}

%% Abstract
\begin{abstract}
Solving \wregulargames is a core problem in formal verification, particularly within model checking and synthesis methods that have several industrial-scale applications. In this paper, we propose a constraint-based approach for games on graphs, built upon a novel propagation algorithm that eliminates opponent cycles from the game graph. This approach yields an efficient method that directly exploits the duality between players in a parity game. Our implementation within a Lazy Clause Generation framework outperforms all existing constraint-based methods for parity games. Furthermore, it performs competitively against the most specialized, non-CP-based algorithms---often matching and sometimes exceeding their performance---while retaining the inherent flexibility of general constraint-based approaches. 

Because new constraints can be seamlessly integrated, our method provides a highly flexible framework for refining existing problem formulations to search for preferred solutions. This extensibility allows the framework to solve multi-conditioned games, such as those combining parity conditions with cost or resource-bounded requirements. Finally, we present the theoretical foundations of our approach, a comprehensive experimental evaluation, and an analysis of different algorithmic variants.
\end{abstract}

%%Graphical abstract
\begin{graphicalabstract}
%\includegraphics{grabs}
\end{graphicalabstract}

%%Research highlights
\begin{highlights}
\item Research highlight 1
\item Research highlight 2
\end{highlights}

%% Keywords
\begin{keyword}
Constraint programming  \sep Parity games \sep Verification.
\end{keyword}

\end{frontmatter}

%% Add \usepackage{lineno} before \begin{document} and uncomment 
%% following line to enable line numbers
%% \linenumbers

%% main text
%%
\newpage
\section{Games on Graphs}
\subsection{Preliminaries}
\subsection{A Constraint-base approach}
\subsection{Objectives / Winning conditions}
\subsubsection{Cualitative conditions}
\begin{itemize}
  \item Buchi / coBuchi : Definition from constraints
  \item Rabin / Street : Definition from constraints
  \item Parity : Definition from constraints
\end{itemize}
\subsubsection{Quantitative conditions}
\begin{itemize}
  \item Energy : Definition from constraints
  \item Mean-payoff : Definition from constraints
\end{itemize}


\section{Constraint method}
\subsection{Boolean version}
\subsection{Integer version}
\subsection{Search strategies / heuristics}

\section{Experments}

%% For citations use: 
%%       \citet{<label>} ==> Lamport (1994)
%%       \citep{<label>} ==> (Lamport, 1994)
%%
% Example citation, See \citet{lamport94}.

%% If you have bib database file and want bibtex to generate the
%% bibitems, please use
%%
%%  \bibliographystyle{elsarticle-harv} 
%%  \bibliography{<your bibdatabase>}

%% else use the following coding to input the bibitems directly in the
%% TeX file.

%% Refer following link for more details about bibliography and citations.
%% https://en.wikibooks.org/wiki/LaTeX/Bibliography_Management

\begin{thebibliography}{00}

%% For authoryear reference style
%% \bibitem[Author(year)]{label}
%% Text of bibliographic item

\bibitem[Lamport(1994)]{lamport94}
  Leslie Lamport,
  \textit{\LaTeX: a document preparation system},
  Addison Wesley, Massachusetts,
  2nd edition,
  1994.

\end{thebibliography}
\end{document}

\endinput
%%
%% End of file `elsarticle-template-harv.tex'.


